\bookchapter{第二十二章：AI如何於手機中工作}

\chapterintro{歡迎至手機最奇之“魔術現場”，主角乃：人工智能。君每日所用之“一拍即得”美照、張口即來之語音助手、越用越懂君之輸入法，背後皆有此隱形大腦於焉。}

\booksection{一、AI相機：君之“後期暗房”已入手機}

今之AI計算攝影，更似將一專業修圖師+攝影師團隊徑塞入快門鍵中。於君按下快門前，AI已始工作：掃描場景，識出“此乃藍天”、“此乃人臉”、“此乃夜景”。

\booksection{二、語音助手：君之“聲音祕書”何以懂君？}

\begin{enumerate}
  \item \textbf{聽清（語音識別）}：將聲波轉爲文字。
  \item \textbf{聽懂（自然語言理解）}：理解意圖，提取關鍵實體。
  \item \textbf{執行與回答}：調用對應App或服務，再以語音合成"說"與君聽。
\end{enumerate}

\booksection{三、端側AI vs. 雲端AI}

端側AI將計算留於手機中，速快、隱私高、可離線。雲端AI需聯網，然能處理更繁雜之任務。

\booksection{四、AI大模型與手機：君之隨身智能助手}

2025-2026年，AI大模型（如DeepSeek R1、ChatGPT-5、文心一言5.0、通義千問3等）之爆發深刻改變矣手機之使用方式。手機不再只為執行指令之工具，而乃能理解、推理與創作之智能夥伴。端側AI芯片（NPU算力突破80 TOPS）使大模型推理首次於手機本地跑通，無需聯網也能用AI。

\booksubsection{何爲AI大模型？}

AI大模型（Large Language Model，LLM）乃經海量文本數據訓練而出之AI系統，能理解並生成自然語言。君可想象其爲一讀過互聯網上幾乎所有文字之“超級學霸”，能答問、撰文、翻譯、編程、析數據。

\booksubsection{手機上之AI大模型應用：}

\begin{itemize}
  \item \textbf{對話助手}：DeepSeek、ChatGPT、豆包、Kimi等App，可隨時提問、翻譯、撰文案、總結文章。比傳統搜索引擎更智——其不只予君鏈接，乃徑予君答案。
  \item \textbf{寫作與創作}：令AI助君撰郵件、總結會議紀要、潤色文章、生成創意文案。
  \item \textbf{圖像生成與編輯}：AI一鍵消除照片中之路人、擴展圖片邊界、風格轉換、甚至依文字描述生成圖片。
  \item \textbf{實時翻譯}：通話實時翻譯、面對面對話翻譯、拍照翻譯，出國旅行不再有語言之礙。
  \item \textbf{學習輔導}：AI可釋概念、出題練習、批改作業，乃隨身家教。
\end{itemize}

\booksubsection{端側大模型：手機自運AI}

最新之旗艦手機已能於本地運行小型AI大模型（通常7B參數左右），所致之變乃：

\begin{itemize}
  \item \textbf{不需網絡}：斷網亦能用AI功能，飛機上、地鐵裏皆可。
  \item \textbf{隱私更安全}：君之對話與數據不會上傳雲端，全於手機本地處理。
  \item \textbf{響應更速}：無網絡延遲，AI之答幾乎即時。
  \item \textbf{當前限制}：端側模型之能尚不及雲端大模型，繁雜任務仍需聯網。然隨芯片AI算力之升，端側模型將愈強。
\end{itemize}

\booksubsection{系統級AI集成：}

\begin{itemize}
  \item \textbf{蘋果Apple Intelligence}：iOS 18引入之AI系統，集成寫作工具、圖像魔法棒、Siri增強、通知摘要等功能。端側處理優先，繁雜任務方上雲端。
  \item \textbf{華爲盤古大模型}：集成於鴻蒙系統，小藝助手得大模型之能，支持文檔摘要、智慧識屏、AI修圖等。
  \item \textbf{小米澎湃AI}：系統級AI之能，支持AI搜圖、AI摘要、AI實時字幕、端側大模型對話等。
  \item \textbf{三星Galaxy AI}：通話實時翻譯、筆記AI摘要、圖片生成編輯、搜索圈選等功能。
\end{itemize}

\booksection{五、數個典型應用場景與實操}

爲令君更直觀感受手機AI之“超能”，下文以五個具體場景，示AI何以於日常生活中“出手相救”：

\booksubsection{場景一：實時翻譯——出國旅行之隨身翻譯官}

\textbf{實際表現}：於巴黎點餐不諳法語？iPhone之“翻譯”App支持面對面雙人對話模式，雙方各觀己方語言之字幕，甚至可同聲傳譯。安卓陣營之三星Galaxy AI通話實時翻譯能令君於電話中徑說中文，對方所聞乃英文，反之亦然。

\textbf{背後原理}：語音識別（ASR） 機器翻譯（NMT） 語音合成（TTS）三步流水線，端到端延遲已可控於1秒內。

\booksubsection{場景二：AI消除——拯救被路人“毀掉”之遊客照}

\textbf{實際表現}：於景點拍一絕美自拍，然背景中擠滿他客？用“魔法消除”圈選路人，手機將自動識別並“塗去”路人，再智能填充背景，效果常令人驚豔。

\textbf{背後原理}：圖像分割（Inpainting）模型先識別並刪除選中區域，而後依周圍像素與語義理解“腦補”出合理之填充內容。此背後之AI模型訓練數據量達數百萬張圖片。

\booksubsection{場景三：AI修圖——一鍵成“攝影大師”}

\textbf{實際表現}：拍了一張陰天灰濛濛之風景照？AI“魔法天空”可一鍵將天空換爲晴空、晚霞甚至星空。AI擴圖可令豎構圖變爲橫構圖，手機“腦補”出畫面外之場景。

\textbf{背後原理}：基於生成式AI（Generative AI）之擴散模型（Diffusion Model），理解圖像之語義內容後生成新像素。此乃當前AI圖像處理之“最前沿”。

\booksubsection{場景四：文檔處理——辦公族之效率倍增器}

\textbf{實際表現}：會議畢，將會議錄音予AI，其能自動整理出結構化之會議紀要、提取待辦事項、生成摘要。一份10頁之英文PDF報告，AI可於10秒內爲君提煉出核心結論。

\textbf{背後原理}：長上下文大模型（Long Context LLM）可“讀完”數萬字之文檔，行多輪推理與信息抽取。Apple Intelligence、華爲小藝、WPS AI、騰訊元寶等皆具此能。

\booksubsection{場景五：AI搜圖與識物——君之“視覺搜索引擎”}

\textbf{實際表現}：於街上見一種不識之花，相機一開，AI立告君“此乃木槿花”。相冊中翻尋去年於海邊所拍之照片，搜索“大海”，AI自動將所有含海之照片尋出，無需手動翻尋。

\textbf{背後原理}：計算機視覺（CV）模型對圖片內容行向量化（Embedding），再與文字查詢行相似度匹配。此等“多模態搜索”乃手機相冊管理之一大飛躍。

\booksection{六、提示詞優化：令AI聽君指揮之“溝通公式”}

AI答之質，大抵繫於君何以“問”。同一之問，問法異，則AI所答天壤之別。下文吾輩介紹一套\textbf{最實用、最易上手之“角色-背景-任務-限制”溝通思路}，並配以大量實例。

\begin{tcbinfo}{何故需框架？}
多人用AI如用百度——敲數關鍵詞入之。然AI非搜索引擎，其乃一“閱讀理解後回答”之語言模型。\textbf{予其更多上下文、說清君欲何、限其勿跑偏}，其即能予君精準之結果。框架乃將君腦中“好答案長何樣”譯爲AI能解之指令。

\end{tcbinfo}

\booksubsection{角色（Role）：告AI其爲何人}

予AI一具體之“身份”，其即調用對應領域之專業知識與表達方式。\textbf{角色設定愈具體，輸出愈專業}。

\booksubsection{實用示例：}

\begin{itemize}
  \item "君乃一位\textbf{有10年經驗之小米手機維修工程師}。"（技術問題）
  \item "君乃一位\textbf{擅長將繁雜概念講簡之中學物理老師}。"（科普問題）
  \item "君乃一位\textbf{審美好、閱書多、自有風格之圖書編輯}。"（寫作問題）
  \item "君乃一位\textbf{每日爲上班族做早餐之烹飪博主}。"（生活問題）
  \item "君乃一位\textbf{耐心、溫和、永不哂新手之中醫}。"（健康問題）
\end{itemize}

\textbf{何以有效？}角色設定了AI之\textbf{知識範圍}、\textbf{表達風格}與\textbf{專業深度}。無角色之AI乃一“全才然啥皆不精”之通才；有角色，其能“入”彼身份，如那位專業人士般思考。

\booksubsection{背景（Background）：告AI君之情況}

將“君爲何人、君之現狀、君之約束條件”先說清楚，AI方能予\textbf{真適合君}之答案，而非泛泛而談。

\booksubsection{實用示例：}

\begin{itemize}
  \item "吾之手機乃\textbf{小米14，購半年，近玩遊戲時發燙嚴重}。"（設備情況）
  \item "吾乃\textbf{一60歲之退休婦人，目力不佳，不善用智能手機}。"（用戶畫像）
  \item "吾正\textbf{備2026年6月之英語六級考試，基礎一般}。"（學習背景）
  \item "吾開了一家\textbf{三線城市之小咖啡店，月流水3萬，僱兩名員工}。"（商業背景）
  \item "吾\textbf{甫誕嬰在家坐月子，睡眠嚴重不足}。"（生活階段）
\end{itemize}

\textbf{何以有效？}背景信息令AI知\textbf{“爲誰服務”}與\textbf{“於何條件下”}。同一題，予學生與予家長講法不同；予城人與予鄉人建議亦異。背景缺，則AI只能予“最大公約數”之答案，常不痛不癢。

\booksubsection{任務（Task）：明告AI做何事}

此步當以\textbf{動詞開頭}（寫、列、對比、翻譯、解釋、總結……），將任務描述得\textbf{具體、可執行}。含糊之任務唯換來含糊之答案。

\booksubsection{實用示例：}

\begin{itemize}
  \item "請\textbf{列出5個吾應立即關閉之權限}，並說明緣由。"（動作明確）
  \item "請\textbf{將下段英文郵件譯爲中文}，保持商務口吻。"（任務具體）
  \item "請\textbf{對比此四款3000元價位之手機}，予出一張優缺點對比表。"（格式明確）
  \item "請\textbf{以小學生能懂之言解釋‘何爲5G’}。"（風格明確）
  \item "請\textbf{爲吾之早餐店寫三條短視頻文案}，分別針對上班族、情侶、帶孩子之家長。"（細分明確）
\end{itemize}

\textbf{何以有效？}任務部分乃AI行動之“指令中樞”。\textbf{動詞愈清晰、要求愈具體，AI愈知“交付何物”}。模糊之“稍弄”、“搞個”、“看看”乃AI大忌——其將揣測君意，結果常跑偏。

\booksubsection{限制（Constraints）：告AI勿做何事}

“限制”聽之似“約束”，實乃\textbf{令AI收斂至最適合君之輸出}之利器。常見維度有：字數、格式、語氣、避免項、必含項。

\booksubsection{實用示例：}

\begin{itemize}
  \item "\textbf{不超過200字，分點列出}。"（長度+格式）
  \item "\textbf{語氣當親切自然，如與友聊天}，勿用‘先、次之、再次’此等官話。"（語氣）
  \item "\textbf{勿用專業術語}，遇專業名詞當順便釋之。"（用詞限制）
  \item "\textbf{每段不超過3行}，方便手機閱讀。"（排版）
  \item "\textbf{必含：使用場景、優缺點、與吾預算之匹配度}。"（必含項）
  \item "\textbf{勿予具體投資建議}，只釋原理。"（邊界）
\end{itemize}

\textbf{何以有效？}無限制，AI易“自由發揮”——或太學術、或太長、或跑題。明確之限制=明確之“交付標準”，雙方對齊預期，結果方可控。

\booksubsection{提示詞質量評估：好提示詞之5個特徵}

寫好提示詞非一次性之事，\textbf{寫完先自查此五條}，再發與AI：

\begin{enumerate}
  \item \textbf{具體而非抽象}：非"寫篇文章"，乃"寫一篇800字、3個小節、面向新手媽媽之手機選購指南"。
  \item \textbf{含背景而非孤立}：告AI"吾爲何人、何故問此"，其方能予"爲君定製"之答案。
  \item \textbf{有明確輸出格式}：爲表格、列表、段落、抑或JSON？先說之。
  \item \textbf{限制到位}：字數、風格、避項，三選一即有明顯提升。
  \item \textbf{可被AI驗證}：問畢，令AI複述一遍君之要求；說不清楚即未說清楚。
\end{enumerate}

\booksubsection{常見問題與解法}

\begin{center}
\begin{tabularx}{\linewidth}{|X|X|}
\hline
\textbf{症狀} & \textbf{解法} \tabularnewline
\hline
回答太泛、不實用 & 加背景（君爲何人/何場景） \tabularnewline
回答太學術/太官方 & 加語氣限制 + 拒"官話"清單 \tabularnewline
回答跑題/答非所問 & 以動詞重述任務，並要求"只答此" \tabularnewline
回答太長看不完 & 硬性限字數（200字內） \tabularnewline
AI一本正經胡說 & 加"若不確定請標‘吾不確定’" + 重要事實自查 \tabularnewline
第一版不滿意 & 勿重寫，於原提示詞後追加"再…一些" \tabularnewline
\hline
\end{tabularx}
\end{center}

\booksubsection{三個完整場景案例}

\begin{tcbinfo}{場景一：生活助手（早餐搭配建議）}
\textbf{提示詞}：

君乃一位擅長營養搭配之早餐博主（\textbf{角色}）。吾乃一位上班族媽媽，每日晨只有15分鐘做早餐，家有6歲之孩子與60歲之婆婆（\textbf{背景}）。請\textbf{列出7天不重樣之早餐搭配}（\textbf{任務}）。每份早餐要求：(1) 準備時間不超過15分鐘 (2) 營養均衡（蛋白+碳水+蔬果）(3) 老少皆宜 (4) 食材於普通超市能買到（\textbf{限制}）。

\textbf{何以有效}：角色定專業度，背景予三類人共性，任務明交付物，4條限制保可執行。

\end{tcbinfo}

\begin{tcbinfo}{場景二：學習輔助（英語六級備考）}
\textbf{提示詞}：

君乃一位\textbf{經驗豐富之英語六級老師}（角色）。吾乃一普通本科生，\textbf{四級480分，六級目標500，尚餘3月}（背景）。請\textbf{爲吾制定一份可執行之12周備考計劃}（任務）。要求：(1) 每週明聽說讀寫之時間分配 (2) 予出每月之重點突破方向 (3) 推薦3本合適之複習資料（中文說明緣由） (4) 表格形式輸出（限制）。

\textbf{何以有效}：角色定專業背景，背景予出分數與時間錨點，任務明“分階段計劃”，限制要求“表格化”與“可執行”。

\end{tcbinfo}

\begin{tcbinfo}{場景三：工作效率（會議紀要整理）}
\textbf{提示詞}：

君乃一位\textbf{資深行政助理}（角色）。吾剛開完一個1小時之部門會議，\textbf{討論Q2營銷推廣方案}（背景）。\textbf{請將以下會議記錄整理成結構化紀要}（任務）。要求：(1) 分“決議事項 / 待辦事項 / 風險與異議”三部分 (2) 待辦事項標出負責人與截止時間（若未提及則寫“待確認”） (3) 總字數控於600字內 (4) 以markdown格式輸出（限制）。

\textbf{何以有效}：角色定行政視角，背景予會議主題，任務明結構化輸出，4條限制令結果徑可複用。

\end{tcbinfo}

\textbf{進階玩法：令AI爲君之“私人助理”}\par

試此等進階用法，以體驗AI之真正實力：

\begin{itemize}
  \item \textbf{多輪對話優化}：對AI之回答不滿意時，可曰"此太專業矣，請以小學生能懂之言再釋一次"。
  \item \textbf{鏈式思考（CoT）}：令AI"一步步思考"，回答質量通常更高："請一步步析此數學題"。
  \item \textbf{多AI協作}：以A寫初稿，B潤色，C翻譯——組合使用效果更佳。
  \item \textbf{系統提示詞}：於AI設置中預設君之身份、偏好、忌諱等，令回答更合君之需。
\end{itemize}

\booksection{七、AI之侷限與挑戰：保持清醒之判斷力}

AI雖強，然絕非“無所不能”。清醒識其侷限，方能更好用之：

\begin{itemize}
  \item \textbf{"幻覺"問題}：AI或一本正經地胡言（專業術語曰Hallucination）。其將編造不存在之歷史事件、錯誤之引言、虛假之統計數據。重要信息必須人工覈實。
  \item \textbf{知識截止}：多數AI之訓練數據有截止日期（如2024年初），其不知之後發生之事。
  \item \textbf{推理能力有限}：繁雜之多步推理、數學證明、邏輯陷阱，AI仍易出錯。
  \item \textbf{隱私風險}：於雲端AI服務中，君之對話或被用於模型訓練。敏感信息當優先用端側AI。
  \item \textbf{偏見與價值觀}：AI或帶訓練數據中之偏見，需保持獨立思考。
\end{itemize}

\begin{tcbinfo}{AI使用小貼士}
\begin{itemize}
  \item AI生成之內容或不準（"幻覺"問題），重要信息務必覈實。
  \item 勿向AI輸入個人隱私信息（身份證號、密碼、銀行信息等）。
  \item 善用提示詞（Prompt）：愈具體、愈清晰之描述，AI之回答質量愈高。
  \item AI乃工具非權威，保持獨立思考與判斷力。
  \item 遇關鍵決策（醫療、法律、財務），諮詢專業人士而非依賴AI。
\end{itemize}

\end{tcbinfo}

\begin{tcbpractice}{來挑戰：三道AI實戰題}
\begin{enumerate}
  \item \textbf{挑戰一：複雜指令}：對語音助手曰："明晨7點提醒吾帶傘，並將今日午後3點之日曆改至4點。" 觀其能否一次完成多步操作。
  \item \textbf{挑戰二：AI修圖}：拍一張有路人之遊客照，用"魔法消除"或"AI橡皮擦"功能，將路人清除乾淨。觀AI對背景之"腦補"效果。
  \item \textbf{挑戰三：提示詞優化}：先以簡單提示詞"寫一段旅遊攻略"問AI，再以CRISPE框架寫一詳細提示詞。較兩次回答之質量差異。
\end{enumerate}

\end{tcbpractice}

\begin{tcbinfo}{案例 2：換機數據遷移全程}
\begin{itemize}
  \item \textbf{場景}：iPhone 13 小米 15 跨品牌換機。
  \item \textbf{步驟}：(1) 舊機 iCloud/微信自帶備份 (2) 退 Apple ID + 閉"尋" (3) 新機用"小米換機"App 掃碼 (4) 擇遷移項 (5) 微信單獨用自帶"聊錄遷移" (6) 驗鍵數據（通訊錄/相冊/微信） (7) 舊機加密擦。
  \item \textbf{預期果}：數據零丟失，舊機可安全轉賣。
\end{itemize}

\end{tcbinfo}
